\begin{center}
\SourceLeaf{025b}\pdfbookmark[2]{7.}{article-007}7.
\end{center}
In perihelio itaque anomalia vera, anomalia excentrica et anomalia media sunt $= 0$;
crescente dein vera, etiam excentrica et media augentur, ita tamen, ut excentrica minor
maneat quam vera, mediaque minor quam excentrica, usque ad aphelium, ubi omnes tres
simul fiunt $= 180^0$; hinc vero usque ad perihelium excentrica perpetuo est maior quam
vera, mediaque maior quam excentrica, donec in perihelio omnes tres fiant $= 360^0$, sive,
quod eodem redit, omnes iterum $= 0$. Generaliter vero patet, si anomaliae verae $v$
respondeat excentrica $E$ mediaque $M$, verae $360^0-v$ respondere excentricam $360^0-E$
atque mediam $360^0-M$. Differentia inter anomaliam veram et mediam $v-M$ \emph{aequatio
centri} appellatur, quae itaque a perihelio ad aphelium positiva, ab aphelio ad
perihelium negativa est, in perihelio ipso autem et aphelio evanescit. Quum igitur $v$ et
$M$ circulum integrum a $0$ usque ad $360^0$ eodem tempore percurrant, tempus revolutionis
unius, quod et \emph{tempus periodicum} dicitur, in diebus expressum invenitur, dividendo
$360^0$ per motum diurnum $\frac{k\BareRoot{(1+\mu)}}{a^{\frac{3}{2}}}$, unde patet, pro
corporibus diversis circa Solem revolventibus quadrata temporum periodicorum cubis
distantiarum mediarum proportionalia esse, quatenus ipsorum massas, aut potius massarum
inaequalitatem, negligere liceat.
